\section{Guidance}
\subsection{Research}

\paragraph{Measuring devices}
\begin{itemize}
    \item strain-gauge
    \item accelerometer
\end{itemize}

\paragraph{Vibration motor}
\begin{itemize}
    \item generator for the waves (linear vs circular?) possibly build? 555 timer + amp + speaker?
    \item Research: vibro-tacticle feedback paper. forces and frequency being used?
\end{itemize}

\subsection{Writing the outline}
\paragraph{Motivation}
Start with clinical problem and move in to the more focused. By the time they're done with reading my intro, they agree with why I'm making this solution

Amputations -> Loss of sensations -> How can we regain that sensation? One of the best tools is using the prosthetic itself, and we can use vibrotactile feedback

describe itap, amputations, etc
ITAP as the vehicle for vibrotactile feedback. Vibrotactile feedback is one of the solutions to the problem of loss of "senses"

\paragraph{Risks}
\begin{itemize}
    \item Soldering risks (burning)
    \item Lone-working
    \item Long-hours
\end{itemize}